\section{The exact contraction ratio for Euclidean configurations}
\label{sec:euclidean-contraction}

The flat hypothesis and the contracting hypothesis cannot be separated
merely by testing Euclidean distance arrays on the matrix labels.
In fact those arrays satisfy a substantially stronger contraction
than the threshold in Section~\ref{sec:alternative-audit}.
We compute the sharp constant directly by finite Fourier analysis.
The coefficient field of tensors plays no role in this calculation:
the matrices below are labels over $\F_2$, while distances are real.

\begin{proposition}[Sharp Euclidean matrix contraction]
\label{prop:euclidean-matrix-contraction}
Let $\mu\geq1$, let $G=\operatorname{Mat}_{\mu}(\F_2)$,
and let $f:G\to H$ take values in a real Hilbert space.  Set
\[
 D_f(r)^2=
 \mathbb E_{a\in G,\ \operatorname{rk}W=r}
          \|f(a+W)-f(a)\|^2.
\]
Then
\[
 D_f(\mu)^2\leq2(1-2^{-\mu})D_f(1)^2.
\]
The constant is sharp, both for arbitrary maps and for injective
maps into finite-dimensional Euclidean space. For $\mu=3$, in particular,
$D_f(3)\leq\sqrt{7/4}\,D_f(1)$.
\end{proposition}

\begin{proof}
The finite image of $f$ lies in a finite-dimensional affine
subspace, so ordinary vector-valued Fourier expansion suffices.
For $Z\in G$ put
\[
 \chi_Z(a)=(-1)^{\operatorname{tr}(Z^{\mathsf T}a)},\qquad
 \widehat f(Z)=\mathbb E_a\chi_Z(a)f(a),\qquad
 \lambda_r(Z)=\mathbb E_{\operatorname{rk}W=r}\chi_Z(W).
\]
The trace in the exponent is taken in $\F_2$.
Orthogonality of these real characters and Parseval's identity give
\begin{equation}\label{eq:euclidean-fourier-increments}
 D_f(r)^2=
 2\sum_{Z\in G}\|\widehat f(Z)\|^2(1-\lambda_r(Z)).
\end{equation}
The term $Z=0$ contributes zero.

Write $Q=2^\mu$ and suppose $Z\ne0$ has rank $s$.
Every nonzero rank-one matrix over $\F_2$ has a unique expression
$uv^{\mathsf T}$ with $u,v\ne0$.
There are $2^{\mu-s}-1$ nonzero $u$ for which
$Z^{\mathsf T}u=0$.  For those $u$, the character sum over
$v\ne0$ is $Q-1$; for every other $u$ it is $-1$.
Consequently
\[
 \lambda_1(Z)=
 \frac{Q^2 2^{-s}-2Q+1}{(Q-1)^2},\qquad
 1-\lambda_1(Z)=
 \frac{Q^2(1-2^{-s})}{(Q-1)^2}
 \geq\frac{Q^2}{2(Q-1)^2}.
\]

We next show
\[
 |\lambda_\mu(Z)|\leq\frac1{Q-1}.
\]
Choose a nonzero column $v$ of $Z$ and condition a uniformly
random invertible $W$ on all columns except the corresponding
one.  The fixed columns span a hyperplane $U$, and the remaining
column is uniform on $\F_2^\mu\setminus U$.
If the character $w\mapsto(-1)^{v\cdot w}$ is nontrivial on
$U$, its average on that complement is zero.
If $U=v^\perp$, its average there is $-1$.
The other columns contribute a factor of absolute value one.
The hyperplane $U$ is uniform among the $Q-1$ hyperplanes,
because each has the same number of ordered bases.  The bound
follows.  This reasoning also includes $\mu=1$, when $U=\{0\}$.

It follows for every nonzero $Z$ that
\[
 \frac{1-\lambda_\mu(Z)}{1-\lambda_1(Z)}
 \leq
 \frac{Q/(Q-1)}{Q^2/(2(Q-1)^2)}
 =2\frac{Q-1}{Q}.
\]
Multiplying these inequalities by the nonnegative Fourier
weights in \eqref{eq:euclidean-fourier-increments} and summing
proves the claimed contraction.  This also handles $D_f(1)=0$
without dividing by that quantity.

For sharpness, take a rank-one $Z=uv^{\mathsf T}$.
For uniform invertible $W$, the vector $Wv$ is uniform among
the nonzero vectors, so $\lambda_\mu(Z)=-1/(Q-1)$.
The displayed lower bound on $1-\lambda_1(Z)$ is an equality
when $s=1$.  Thus $f(a)=\chi_Z(a)$ attains the ratio.
For an injective example, take all the matrix-unit characters:
\[
 f(a)=\bigl(\chi_{E_{ij}}(a)\bigr)_{1\leq i,j\leq\mu}
       \in\mathbb R^{\mu^2}.
\]
This map records every binary matrix entry and is injective.
Every nonconstant Fourier component has rank-one index,
so the same ratio is attained exactly.
\end{proof}

For every $\mu\geq3$ the sharp Euclidean ratio
$\sqrt{2(1-2^{-\mu})}$ is strictly smaller than $\sqrt\mu$.
Thus all finite Euclidean metrics are compatible with a
contraction inequality of the required strict kind on these
labels.  This does not prove that any field satisfies both
the flat and contracting hypotheses: contraction quantifies
over every price, which need not be a Euclidean metric.
It shows exactly why the presence of all Euclidean prices,
by itself, does not refute contraction.

\subsection{The full curvature threshold is attained in a building}

The stronger Hilbert bound does not extend to all CAT(0) targets.
The following classical construction identifies the precise obstruction.
Its nonpositive-curvature input is the theorem of
Haettel--Kielak--Schwer that the orthoscheme complex of a bounded,
graded, complemented modular lattice is CAT(0)
\cite[Theorem~5.6, pp.~25--26]{HKS2016}.
We prove the needed distance formula, including the fact that paths
through other simplices cannot make that distance smaller.

\begin{lemma}[Subspace distances in a CAT(0) complex]
\label{lem:subspace-cat0}
Let $F$ be a finite field and $V$ a finite-dimensional $F$-vector
space. Its subspaces, with distances
\[
 \rho(U,W)^2=\dim U+\dim W-2\dim(U\cap W),
\]
embed isometrically in a complete CAT(0) space.
\end{lemma}

\begin{proof}
The assertion is immediate if $\dim V=0$, so suppose $n=\dim V>0$.
Let $L$ be the lattice of all subspaces of $V$.
It is bounded, with minimum $0$ and maximum $V$, and graded by
vector-space dimension. Meets are intersections and joins are sums.
For $U\subseteq W$ the modular identity
$U+(Z\cap W)=(U+Z)\cap W$ follows directly by writing a vector
in the right-hand side as a sum of a vector in $U$ and one in $Z$.
Every subspace has a complementary subspace, by extension of a basis.
Thus $L$ has all hypotheses of the cited theorem.

For clarity, its orthoscheme complex $X$ has one simplex for each
chain of subspaces. A maximal chain
$0=U_0\subset U_1\subset\cdots\subset U_n=V$
is made isometric to the simplex in $\mathbb R^n$ with vertices
$0,e_1,e_1+e_2,\ldots,e_1+\cdots+e_n$; the metric on $X$ is the
induced length metric \cite[Definition~2.17, p.~6]{HKS2016}.
The theorem makes $X$ CAT(0). Since $F$ is finite, $X$ is a finite
complex of compact Euclidean simplices and is compact, hence complete.

Fix a basis $b_1,\ldots,b_n$ of $V$, and let $Q$ be the subcomplex
of coordinate subspaces for this basis. Under
$\operatorname{span}\{b_i:i\in S\}\mapsto\mathbf1_S$, its chain
simplices triangulate the Euclidean cube $[0,1]^n$.
Put $F_i=\operatorname{span}(b_1,\ldots,b_i)$ and define a subset
$p(U)\subseteq\{1,\ldots,n\}$ by
\[
 i\in p(U)\quad\Longleftrightarrow\quad
 \dim(U\cap F_i)-\dim(U\cap F_{i-1})=1.
\]
Each dimension increment is zero or one. If $U\subseteq W$, the
natural map
\[
 (U\cap F_i)/(U\cap F_{i-1})
       \longrightarrow (W\cap F_i)/(W\cap F_{i-1})
\]
is injective; consequently $p(U)\subseteq p(W)$.
Also $|p(U)|=\dim U$, and $p$ fixes every coordinate subspace.
Extending this vertex map affinely on chain simplices defines a
retraction $p:X\to Q$. It is an isometry on each simplex, because
it preserves all vertex ranks and their differences. It is therefore
nonexpanding for the length metrics: apply it to a path, subdivide
that path among simplices, and take the infimum of lengths.
It follows that distances between points of $Q$ computed in $X$
are their Euclidean cube distances. Indeed the retraction gives
one inequality, and a straight segment in $Q$ gives the other.

Given any two subspaces $U,W$, choose a basis of $U\cap W$,
extend it separately to bases of $U$ and $W$, and extend their
union to a basis of $V$. Then $U,W$ belong to the corresponding
coordinate cube. Their indicator vectors differ in exactly
$\dim U+\dim W-2\dim(U\cap W)$ coordinates. The formula follows.
\end{proof}

\begin{proposition}[Sharp CAT(0) matrix contraction]
\label{prop:cat0-matrix-extremizer}
For every $\mu\ge1$, the metric
\[
 d(A,B)=\sqrt{\operatorname{rk}(A-B)}
       \qquad\bigl(A,B\in\operatorname{Mat}_{\mu}(\F_2)\bigr)
\]
embeds isometrically in a complete CAT(0) space.
It satisfies $D(r)^2=r$ for every $0\le r\le\mu$.
In particular the bound $D(\mu)\le\sqrt\mu\,D(1)$ for CAT(0)
targets is sharp, including when $\mu\ge3$.
\end{proposition}

\begin{proof}
In $V=\F_2^\mu\oplus\F_2^\mu$ associate to $A$ its graph
\[
 U_A=\{(x,Ax):x\in\F_2^\mu\}.
\]
It has dimension $\mu$, while
$\dim(U_A\cap U_B)=\dim\ker(A-B)=\mu-\operatorname{rk}(A-B)$.
Lemma~\ref{lem:subspace-cat0} therefore gives
$\rho(U_A,U_B)^2=2\operatorname{rk}(A-B)$.
Rescale the ambient CAT(0) metric by $1/\sqrt2$.
The displayed formula for $D(r)$ is immediate.
The universal upper bound follows by averaging the squared
quadrilateral inequality on transverse increments, as in
Section~\ref{sec:alternative-audit}.
\end{proof}

Thus CAT(0) embeddability by itself permits the exact extremal
matrix metric that appears in the compactness alternative.
Neither the construction nor the CAT(0) theorem makes this metric
a tensor price. A strict contraction argument must use an
additional restriction on prices beyond their embedding in a
complete CAT(0) space.

\subsection{An exact price test for contraction}

For the tensor price cone, the extremal rank metric gives a finite
and exact test of the contracting hypothesis. This uses the equality
classification already proved in Section~\ref{sec:alternative-audit},
not a claim that all CAT(0) metrics are prices.

\begin{proposition}[The rank-metric gate for contraction]
\label{prop:rank-price-contraction-gate}
Assume the price-space theorem, namely Proposition~C.3 of
\cite{Anthropic2026}, and its CAT(0) equality consequences proved in
Section~\ref{sec:alternative-audit}.
Fix a coefficient field $\kk$ and $\mu\ge3$. The maximum
\[
 R_\mu(\kk)=
 \max\{D_d(\mu)^2:d\text{ is a price over }\kk,\ D_d(1)=1\}
\]
exists and satisfies
\[
 2(1-2^{-\mu})\le R_\mu(\kk)\le\mu.
\]
Moreover $R_\mu(\kk)=\mu$ if and only if the metric
$d(A,B)=\sqrt{\operatorname{rk}(A-B)}$ is a price over $\kk$.
Consequently the contracting hypothesis holds exactly when this
rank metric is a nonprice for at least one integer $\mu\ge3$.
\end{proposition}

\begin{proof}
The normalized price arrays form a closed set. The diameter bound
in Section~\ref{sec:alternative-audit}, namely
$\mu\sqrt{2^{\mu^2}(2^\mu-1)^2}$, makes it compact.
It is nonempty: the distance of a nonconstant scalar Walsh
character is twice a cut metric, and cut metrics are wire prices
over every coefficient field. Taking a rank-one character and
using Proposition~\ref{prop:euclidean-matrix-contraction} gives
the displayed lower bound after normalization. The averaged
quadrilateral inequality gives the upper bound. Continuity of
$D_d(\mu)^2$ proves attainment.

If the maximum is $\mu$, choose an attaining price. The fixed-$\mu$
equality argument in Section~\ref{sec:alternative-audit} first
forces equality on every transverse quadrilateral and then shows
$d(A,B)^2=\operatorname{rk}(A-B)$ for every pair of matrices.
Conversely that rank metric has $D(1)=1$ and $D(\mu)^2=\mu$.
This proves the equivalence for each $\mu\ge3$.

A nonprice rank metric for such a $\mu$ gives
$1\le\sqrt{R_\mu(\kk)}<\sqrt\mu$, a valid uniform contraction
constant. Conversely, a contraction at a parameter $\mu\ge3$
excludes its rank metric directly. If the given parameter is
$\mu=2$, its rank metric is also excluded directly. The rank metric
on $\operatorname{Mat}_3(\F_2)$ restricts to that on
$\operatorname{Mat}_2(\F_2)$ by adjoining a zero row and column.
Since restrictions of prices are prices, the rank metric for
$\mu=3$ must then be a nonprice as well.
\end{proof}

The Euclidean cross-polytope tests characterize failure of flatness,
whereas Proposition~\ref{prop:rank-price-contraction-gate} tests
contraction using a non-Euclidean metric in a building.
The implication from a nonprice rank metric to a nonprice Euclidean
configuration remains unresolved here. In particular this exact
test does not turn the known implication from nonflatness to
contraction into an equivalence.

\subsection{A common centre and simultaneous spectral restrictions}

A hypothetical rank-metric price has more structure than its individual
Euclidean rectangles. All its diameters have the same midpoint. The
following assertion is intrinsic to the finite distance array; it does
not depend on choosing the building realization above.

\begin{lemma}[The common centre of the rank metric]
\label{lem:rank-metric-common-centre}
Let $\mu\ge1$. In every CAT(0) realization of
$\operatorname{Mat}_{\mu}(\F_2)$ with distances
$d(A,B)=\sqrt{\operatorname{rk}(A-B)}$, there is a unique point $o$
that is the midpoint of every pair with invertible difference.
For every matrix $A$,
\[
 d(o,A)=\frac{\sqrt\mu}{2}.
\]
\end{lemma}

\begin{proof}
The case $\mu=1$ is a single segment. Suppose $\mu\ge2$, fix an
invertible matrix $U$, and write $m_U(A)$ for the midpoint of
$A,A+U$. If $P$ is a rank-one idempotent, then $I+P$ has rank
$\mu-1$. The four vertices
\[
 A,\quad A+UP,\quad A+U,\quad A+U+UP
\]
have successive squared side lengths $1,\mu-1,1,\mu-1$ and both
squared diagonals equal to $\mu$. Applying the CAT(0) midpoint
inequality to each diagonal gives
\[
 d\bigl(m_U(A),m_U(A+UP)\bigr)^2
 \leq\frac14\left(2+2(\mu-1)-2\mu\right)=0.
\]
Indeed, first bound the squared distances from the midpoint of one
diagonal to each endpoint of the other, then apply the midpoint
inequality on that second diagonal.

Rank-one idempotents span the additive group of all binary matrices:
they include $E_{ii}$ and $E_{ii}+E_{ij}$ for $i\ne j$, whose sums
give every off-diagonal matrix unit. Thus the increments $UP$ span
that group as well. The preceding equality shows that $m_U(A)$ is
independent of $A$; call it $o_U$. It is at distance $\sqrt\mu/2$
from every matrix label.

For another invertible $V$, apply the midpoint inequality to $o_V$
and the segment $A,A+U$. Both endpoint distances are $\sqrt\mu/2$,
so $d(o_U,o_V)^2\leq\mu/4-\mu/4=0$. All the points $o_U$ therefore
agree. Uniqueness follows from uniqueness of the midpoint of any
one full-rank pair.
\end{proof}

\begin{proposition}[All full-rank diameters have fixed tensor values]
\label{prop:rank-centre-spectrum}
Assume the price-space theorem and the full normalized extension
of Proposition~\ref{prop:full-normalized-extension}. Suppose an
ordinary normalized point $\psi$ on
$G=\operatorname{Mat}_{\mu}(\F_2)$ has
\[
 d_\psi(A,B)=\lambda\sqrt{\operatorname{rk}(A-B)},\qquad\lambda>0.
\]
It extends to an ordinary normalized point $\psi'$ on $G\cup\{o\}$
which preserves every old tensor value and satisfies
\[
 d_{\psi'}(o,A)=g/2\quad(A\in G),\qquad g=\lambda\sqrt\mu.
\]
For every $A,B$ with invertible difference and every tensor $T$
supported on $(A,B,o)$, the collinear edge classification gives
\begin{equation}\label{eq:rank-centre-edge-values}
 \psi'(T)=\left(\zeta^{(1/2,1/2,0)}(T)\right)^g.
\end{equation}
With $n=2^{\mu^2}$ one has $g\ge2/n$. For $\mu\ge2$ one also has
\[
 g\le\frac{\log_2 7}{3}<1.
\]
For $\mu=1$, ordinary normalization gives $g=1$.
\end{proposition}

\begin{proof}
Adjoin the midpoint of one full-rank pair using two equally weighted
retractions in Proposition~\ref{prop:full-normalized-extension}.
All old tensor values are preserved. The two new endpoint distances
are at most half their old distance, and the triangle inequality
forces equality. In a CAT(0) realization of the extended price,
Lemma~\ref{lem:rank-metric-common-centre}, rescaled by $\lambda$,
identifies the new point with the common centre. The full three-party
values now follow from Proposition~\ref{prop:collinear-edge-point}.

The unit tensor on $G\cup\{o\}$ is a restriction of the product of
the $n$ binary wires joining the centre to the matrix labels. Its
normalized exponent is one, so $1\le ng/2$.

For $\mu\ge2$, represent $\F_{2^\mu}$ by multiplication matrices
on its underlying $\mu$-dimensional binary vector space. If
$\alpha\notin\{0,1\}$, its multiplication matrix $W$ and $I+W$
are invertible. Thus $0,I,W$ form a triangle with all three wire
exponents equal to $g$. The product of its three binary wires is
the $2\times2$ matrix-multiplication tensor. Strassen's seven-term
integral decomposition \cite[pp.~354--355]{Strassen1969} works
after reduction in every characteristic. Cover subadditivity
therefore gives $2^{3g}\le7$. Finally, for $\mu=1$ the original
normalized two-party functional is matrix rank, so $g=1$.
\end{proof}

Equation~\eqref{eq:rank-centre-edge-values} fixes full tensor
functionals on every triple associated with an invertible matrix
difference. These are simultaneous necessary conditions on one
spectral point. No incompatibility among them and the remaining
wire values has been established here, so the contraction test
of Proposition~\ref{prop:rank-price-contraction-gate} remains open.

\subsection{Convexity of full tensor values along geodesics}

The extension theorem also controls tensors whose parties move along
different geodesics at the same time. This gives more information than
the convexity of individual wire distances.

\begin{lemma}[Transport across a zero-distance wire]
\label{lem:zero-wire-identification}
Let $\phi$ be a relaxed spectral point, and suppose
$d_\phi(u,v)=0$. Grouping the tensor factors at $u$ and $v$ at
either of these parties does not change the value of any tensor.
The same holds when finitely many factors are transported between
pairs of parties at distance zero.
\end{lemma}
\begin{proof}
Let the factor to be moved have dimension $r$. Adjoin a wire
$\langle r\rangle_{uv}$. A local equality contraction at its
source, with the other end retained as an output index, moves the
factor to the destination. The reverse operation uses another
copy of the same wire. Each wire has value $r^0=1$, so restriction
monotonicity in the two directions gives equality. Apply this
argument successively to finitely many factors. The argument also
allows several factors to arrive at the same party, where their
indices are grouped by tensor product.
\end{proof}

\begin{proposition}[Geodesic convexity of tensor values]
\label{prop:geodesic-tensor-convexity}
Assume the price-space theorem and
Proposition~\ref{prop:full-normalized-extension}.
Let $\psi$ be an ordinary normalized spectral point on a finite
party set $V$. Let $x,y,z:P\to V$ be maps from another finite
party set, and let $0\le t\le1$. Suppose, for every $i\in P$,
\[
 d_\psi(x(i),z(i))=t\,d_\psi(x(i),y(i)),\qquad
 d_\psi(z(i),y(i))=(1-t)d_\psi(x(i),y(i)).
\]
Then for every nonzero $P$-tensor $T$,
\begin{equation}\label{eq:geodesic-tensor-convexity}
 \psi(zT)\le\psi(xT)^{1-t}\psi(yT)^t.
\end{equation}
Here a map groups all tensor factors with the same image, so
coincident party positions are permitted.
\end{proposition}
\begin{proof}
The cases $t=0,1$ follow from Lemma~\ref{lem:zero-wire-identification}.
For $0<t<1$, adjoin one fresh party $i'$ for every $i\in P$.
Use the two retractions fixing $V$ and sending $i'$ to $x(i)$ or
$y(i)$, with weights $1-t,t$. The extension theorem gives a
normalized point $\psi'$ preserving every old tensor value.
Its new wire exponents satisfy
\[
 d_{\psi'}(i',x(i))\le t\,d_\psi(x(i),y(i)),\qquad
 d_{\psi'}(i',y(i))\le(1-t)d_\psi(x(i),y(i)).
\]
The triangle inequality makes these bounds equalities. In a
CAT(0) realization, both $i'$ and the old party $z(i)$ are
therefore the same time-$t$ point on the unique geodesic from
$x(i)$ to $y(i)$. If the endpoints have distance zero the same
conclusion follows directly from the displayed bounds.
Thus $d_{\psi'}(i',z(i))=0$ for every $i$.

Place $T$ on the fresh parties. The upper bound in the extension
theorem is $\psi(xT)^{1-t}\psi(yT)^t$. Transporting each fresh
factor to $z(i)$ does not change its value by
Lemma~\ref{lem:zero-wire-identification}. The value after transport
is $\psi(zT)$, since all old tensor values are preserved. This
proves the inequality, including all coincidences among the maps.
\end{proof}

\begin{corollary}[Inserting a geodesic point preserves the unit exponent]
\label{cor:unit-geodesic-insertion}
Under the same hypotheses, let $S\subseteq V$ contain $x,y$, and
let $z\in V$ lie on their geodesic. Then
\[
 g_\psi(S\cup\{z\})=g_\psi(S).
\]
The conclusion can be iterated for any finite sequence of such
insertions among the given parties.
\end{corollary}
\begin{proof}
Apply Proposition~\ref{prop:geodesic-tensor-convexity} to a unit
tensor with one party for each element of $S$ and one extra party.
Keep the $S$ parties fixed. Move the extra party from $x$ to $y$
through $z$. At both endpoints, grouping its repeated diagonal
index produces a tensor locally equivalent to the unit tensor
on $S$. The middle tensor is the unit tensor on $S\cup\{z\}$,
up to the same harmless grouping if $z\in S$. The proposition
gives the upper bound on the exponent. Contracting the factor at
$z$ gives the reverse inequality when $z\notin S$; when $z\in S$
the asserted equality is immediate.
\end{proof}

These statements constrain full tensor values and the exponents of
unit tensors. They do not by themselves provide a common spectral
extension for separately specified flats, nor a contradiction to the
rank-metric price prescriptions above.

\subsection{A maximal scale with exact radial tensor values}

Maximizing the normalized scale gives a further selectable constraint.
It fixes every tensor value on each contracted copy, although it does
not determine tensors whose parties occur at different radii.

\begin{lemma}[The intrinsic radial metric]
\label{lem:rank-radial-metric}
In the CAT(0) realization of
Lemma~\ref{lem:rank-metric-common-centre}, let $tA$ denote the
time-$t$ point of the segment from the common centre to $A$,
where $0\le t\le1$. Then
\begin{equation}\label{eq:rank-radial-metric}
 d(tA,sB)^2=\frac\mu4(t-s)^2+ts\operatorname{rk}(A-B).
\end{equation}
\end{lemma}
\begin{proof}
Put $W=B-A$ and $r=\operatorname{rk}W$. If $0<r<\mu$, choose
invertible $L,R$ over $\F_2$ such that
$W=L\operatorname{diag}(I_r,0)R$, and put $U=LR$.
Then $U$ is invertible and $U+W$ has rank $\mu-r$.
The four labels $A,B,A+U,B+U$ bound a flat rectangle:
its squared side lengths are $r,\mu-r,r,\mu-r$ and both
squared diagonals are $\mu$. Its centre is the common centre
already identified. The radial segments to $A,B$ lie in that
rectangle. Their squared lengths are $\mu/4$, and their Euclidean
inner product is $\mu/4-r/2$. Expansion gives the formula.
For $r=0$ there is only one ray. For $r=\mu$ the two rays are
opposite halves of one diameter. The same formula holds in both
cases, including zero radii.
\end{proof}

\begin{theorem}[Selectable exact radial homogeneity]
\label{thm:radial-full-homogeneity}
Assume the price-space theorem and
Proposition~\ref{prop:full-normalized-extension}. Suppose the metric
$\sqrt{\operatorname{rk}(A-B)}$ on
$G=\operatorname{Mat}_{\mu}(\F_2)$ is a price over $\kk$.
Among ordinary normalized points whose wire metrics are positive
multiples of this metric, the diameter exponent has a maximum
$g_*>0$. There is a maximizing point $\psi_*$ which is pointwise
minimal within the maximizing fibre: if another maximizing point
$\eta$ satisfies $\eta(T)\le\psi_*(T)$ for every $G$-tensor,
then $\eta=\psi_*$.

Put $\lambda_*=g_*/\sqrt\mu$. Let $X$ consist of the points
$(t,A)$, $0\le t\le1$, $A\in G$, with every $(0,A)$ identified
to a single point $o$. There is a coherent family of ordinary
normalized spectral points on all finite subsets of $X$ containing
$o$ and the base copy $G=\{(1,A):A\in G\}$. They preserve every
base tensor value of $\psi_*$ and have wire metric
\[
 D((t,A),(s,B))
 =\lambda_*\sqrt{\frac\mu4(t-s)^2+ts\operatorname{rk}(A-B)}.
\]
For every nonzero $G$-tensor $T$ and $0<t\le1$, its value on the
copy $tG$ is exactly
\begin{equation}\label{eq:radial-full-homogeneity}
 \psi(T\text{ placed on }tG)=\psi_*(T)^t.
\end{equation}
At $t=0$ its grouped value is one. In particular, the unit-tensor
exponent on $tG$ is exactly $t$.
\end{theorem}

\begin{proof}
The ordinary normalized spectrum on $G$ is compact. The conditions
that all wire exponents be proportional to the rank metric are
closed if scale zero is temporarily allowed, since the constant is recovered from
$d_\psi(0,I)/\sqrt\mu$. Price attainment supplies a point with
positive scale. The continuous diameter exponent therefore attains
a positive maximum. Its maximizing fibre $K$ is compact and
nonempty.

Order $K$ pointwise on all tensors. A descending chain has a lower
bound in $K$: the closed conditions $\eta\le\phi$, as $\phi$
runs over the chain, have the finite intersection property, because
any finite subchain has a least element. Compactness supplies their
simultaneous solution. Zorn's lemma gives a pointwise-minimal
element $\psi_*$ of $K$. No uniqueness is asserted.

First adjoin the common centre to $\psi_*$ as in
Proposition~\ref{prop:rank-centre-spectrum}, preserving every old
tensor value, and fix one such normalized extension.
Let $J\subset(0,1)$ be a finite set of radii. Adjoin all the parties
$(t,A)$ with $t\in J$, $A\in G$. Use a single random variable
$u$, uniform on $[0,1]$, to define retractions by
\[
 (t,A)\longmapsto
 \begin{cases}A,&u\le t,\\o,&u>t,\end{cases}
\]
fixing the old base and centre. Only finitely many retractions occur,
with probabilities given by the intervals between the radii in $J$.
Apply the full normalized extension theorem to these retractions.

The new party $(t,A)$ has distances at most $t\,d(o,A)$ and
$(1-t)d(o,A)$ from $o$ and $A$, respectively. The triangle
inequality forces equality. In a CAT(0) realization it is thus
the time-$t$ radial point. Lemma~\ref{lem:rank-radial-metric}
fixes all distances, also between different radii.
Because all parties within one copy use the same threshold, the
extension upper bound on every nonzero $G$-tensor $T$ in that copy is
\[
 \eta_t(T)\le\psi_*(T)^t.
\]
When the copy is sent to $o$, its tensor becomes a nonzero vector
at one party and has value one.

Let $h_t$ be the unit-tensor exponent of the restricted relaxed
point $\eta_t$. The upper bound gives $h_t\le t$.
A full-rank wire in that copy has exponent $tg_*>0$, so $h_t>0$.
Normalization makes $\eta_t^{1/h_t}$ an ordinary point on $G$
whose diameter exponent is $tg_*/h_t$. Maximality of $g_*$ gives
$tg_*/h_t\le g_*$, hence $h_t\ge t$. Therefore $h_t=t$.
The normalized point $\eta_t^{1/t}$ lies in $K$ and is pointwise
at most $\psi_*$. Minimality now forces equality on every tensor,
which proves \eqref{eq:radial-full-homogeneity} simultaneously for
all radii in $J$.

For coherence, take the product of the compact normalized spectra
over all finite radius sets, imposing the closed conditions fixing
base values, radial distances, copy values, and restrictions between
the sets. Every finite collection of these conditions is satisfied
by the preceding construction on the union of its finitely many
radius sets and restriction. Such restrictions remain ordinary
normalized, since they contain the fixed base unit tensor of value
two and their own unit tensor has rank two. Compactness gives one
coherent family for all real radii in $[0,1]$. Restrict it to
arbitrary finite subsets containing the base and centre. The radius
zero assertion follows by grouping into a nonzero vector.
\end{proof}

The maximizing and minimal selections are part of the theorem.
An arbitrary extension need not have these exact values. The result
provides compatible homogeneity under the rank-price hypothesis;
it is not a proof that this hypothesis holds or fails.


\subsection{A sixteen-point obstruction to averaging apartments}

\begin{proposition}[A small rank metric is not \(L_1\)]\label{prop:rank-metric-not-lone}
Let \(H\subset\operatorname{Mat}_3(\F_2)\) consist of the sixteen matrices
\[
 A(a,b,c,d)=
 \begin{pmatrix}0&0&a\\0&b&0\\c&0&d\end{pmatrix},
 \qquad a,b,c,d\in\F_2.
\]
The metric \(d(A,B)=\sqrt{\operatorname{rk}(A-B)}\) is not a nonnegative
combination of cut metrics, and hence has no isometric embedding in
\(L_1\). The same holds for every positive multiple of it and for the
full rank metric on \(\operatorname{Mat}_\mu(\F_2)\), \(\mu\geq3\).
This is an obstruction to convex-conic averaging of Euclidean or tree
apartment metrics; it is not a nonprice certificate.
\end{proposition}
\begin{proof}
Put \(\chi(A(a,b,c,d))=(-1)^{a+c}\).  Reordering coordinates gives
\[
 \operatorname{rk}A
 =\operatorname{rk}\begin{pmatrix}0&a\\c&d\end{pmatrix}+b.
\]
For the \(2\times2\) block the zero matrix is unique, and the two
invertible matrices have \(a=c=1\), so both have character \(+1\).
The total nontrivial character sum is zero. Its signed rank-count
polynomial is therefore \(1-3t+2t^2\). The independent bit \(b\)
multiplies this by \(1+t\); the four signed rank counts on \(H\) are
\((1,-2,-1,2)\).
Translation invariance now gives
\begin{equation}\label{eq:rank-negative-type}
 \sum_{A,B\in H}\chi(A)\chi(B)d(A,B)
 =16(-2-\sqrt2+2\sqrt3)>0.
\end{equation}
For exact positivity, \(2\sqrt3>2+\sqrt2\) reduces by squaring to
\(3>2\sqrt2\), and then to \(9>8\).

For any cut metric \(\delta_S\), since \(\sum_A\chi(A)=0\),
\[
 \sum_{A,B\in H}\chi(A)\chi(B)\delta_S(A,B)
 =-2\left(\sum_{A\in S}\chi(A)\right)^2\leq0.
\]
This contradicts \eqref{eq:rank-negative-type} for every nonnegative
cut combination. On a finite point set, every \(L_1\) metric is such
a combination: integrate the threshold cuts of each real coordinate,
noting that only finitely many distinct cuts occur. Positive scaling
preserves the contradiction, and zero-padding embeds \(H\) isometrically
in every larger matrix space.

Finite Euclidean metrics belong to the cut cone by
Proposition~\ref{prop:no-linear-euclidean-obstruction}; finite tree
metrics do by their edge cuts. Their pullbacks and nonnegative
combinations therefore cannot produce this metric.
\end{proof}

A general theorem that all prices have negative type would turn this
witness into a rank-metric nonprice certificate. No such theorem is
proved or assumed here. In particular, the CAT(0) building realization
and this \(L_1\) obstruction are compatible.
